\section{Additional Results}

\subsection{Additional Audio Generation Results}

\subsubsection{Additional Metrics for Main Results on Sound Effect Generation}\label{sec:app_audio_main_sfx}
We present results on objective metrics and additional subjective metrics for the ``sound effect generation'' experiments in \ref{sec:main_res}.
Table~\ref{tab:app_eval_audio_main_ss_sfx} should be compared against \cref{tab:eval_audio_main_ss_sfx}.

\begin{table}[h]
    \centering
    \adjustbox{max width=\textwidth}{%
    \begin{tabular}{lll cc ccc}
    \toprule
    \multirow{4}{*}{Dataset} & \multirow{4}{*}{Baseline} & \multirow{4}{*}{Type} & \multicolumn{2}{c}{Subjective} & \multicolumn{3}{c}{Objective} \\
    \cmidrule(lr){4-5} \cmidrule(lr){6-8}
    & & & \multicolumn{2}{c}{TA Align} & Quality & VA Align & TA Align \\
    & & & Recall (\%) $\uparrow$ & Precision (\%) $\uparrow$ & AQual $\uparrow$ & IB $\uparrow$ & CLAP $\uparrow$ \\
    \midrule
    \multirow{9}{*}{\shortstack[c]{\OursAudioBenchSReal SFX}}
    & Diff-Foley        & V2A  & - & - & 5.68 & 0.28 & 0.36* \\
    & FoleyCraft        & V2A  & - & - & 6.03 & 0.31 & 0.35* \\
    & VTA-LDM           & V2A  & - & - & 5.97 & 0.30 & 0.35* \\
    & Seeing\&Hearing   & V2A  & - & - & 5.21 & 0.38 & 0.33* \\
    & Seeing\&Hearing   & TV2A & \meanci{63.0}{6.7} & \meanci{56.6}{6.5} & 5.70 & 0.34 & 0.47\\
    & PikaLabs          & V2A  & - & - & 6.45 & 0.18 & 0.28* \\
    & PikaLabs          & TV2A & \meanci{64.5}{7.7} & \meanci{62.4}{8.3} & 6.69 & 0.21 & 0.43\\
    & ElevenLabs        & T2A  & \meanci{82.1}{5.9} & \meanci{78.8}{6.6} & 6.53 & 0.24 & 0.54\\
    \cmidrule{2-8}
    & \OursAudio        & TV2A & \meanci{84.7}{4.8} & \meanci{75.5}{5.0} & 6.76 & 0.38 & 0.51\\
    \midrule
    \multirow{9}{*}{\shortstack[c]{\OursAudioBenchSGen SFX}}
    & Diff-Foley        & V2A  & - & - & 5.48 & 0.18 & 0.17*\\
    & FoleyCraft        & V2A  & - & - & 6.03 & 0.24 & 0.25*\\
    & VTA-LDM           & V2A  & - & - & 6.03 & 0.22 & 0.25*\\
    & Seeing\&Hearing   & V2A  & - & - & 5.21 & 0.38 & 0.27*\\
    & Seeing\&Hearing   & TV2A & \meanci{72.3}{3.8} & \meanci{67.2}{3.8} & 5.49 & 0.37 & 0.42 \\
    & PikaLabs          & V2A  & - & - & 6.18 & 0.16 & 0.26*\\
    & PikaLabs          & TV2A & \meanci{63.0}{4.6} & \meanci{68.7}{5.1} & 6.20 & 0.18 & 0.41 \\
    & ElevenLabs        & T2A  & \meanci{71.0}{3.8} & \meanci{68.7}{3.9} & 6.39 & 0.19 & 0.45 \\
    \cmidrule{2-8}
    & \OursAudio        & TV2A & \meanci{73.9}{3.7} & \meanci{71.8}{3.8} & 6.47 & 0.34 & 0.46 \\
    \bottomrule
    \end{tabular}}
    \caption{
        \textbf{Additional sound effect generation subjective and objective evaluations.} This table reports additional metrics on the \OursAudioBenchSReal and \OursAudioBenchSGen SFX benchmarks. TA and VA are abbreviations for text-audio and video-audio respectively. We put ``*'' on CLAP results when text is not used at inference.
    }
    \label{tab:app_eval_audio_main_ss_sfx}
\end{table}

\begin{table}[h]
    \centering
    \adjustbox{max width=\textwidth}{%
    \begin{tabular}{lll ccc}
    \toprule
    \multirow{4}{*}{Dataset} & \multirow{4}{*}{Baseline} & \multirow{4}{*}{Type} & \multicolumn{3}{c}{Objective} \\
    \cmidrule(lr){4-6}
    & & & Quality & VA Align & TA Align \\
    & & & AQual $\uparrow$ & IB $\uparrow$ & CLAP $\uparrow$ \\
    \midrule
    \multirow{7}{*}{\shortstack[c]{\OursAudioBench SFX}}
    & Diff-Foley        & V2A  & 5.26 & 0.21 & 0.19*\\
    & FoleyCraft        & V2A  & 5.73 & 0.25 & 0.24*\\
    & VTA-LDM           & V2A  & 5.74 & 0.22 & 0.17*\\
    & Seeing\&Hearing   & V2A  & 5.21 & 0.37 & 0.25*\\
    & Seeing\&Hearing   & TV2A & 5.52 & 0.33 & 0.43 \\
    & ElevenLabs        & T2A  & 6.17 & 0.20 & 0.47 \\
    \cmidrule{2-6}
    & \OursAudio        & TV2A & 6.39 & 0.36 & 0.45 \\
    \bottomrule
    \end{tabular}}
    \caption{
        \textbf{Sound effect generation objective evaluation on \OursAudioBench.} TA and VA are abbreviations for text-audio and video-audio respectively. We put ``*'' on CLAP results when text is not used at inference. Similar trends as in \cref{tab:eval_audio_main_ss_sfx,tab:app_eval_audio_main_ss_sfx} are observed.
    }
    \label{tab:app_eval_audio_main_ss_sfx_audiobench}
\end{table}


We first note that on text-audio alignment, \OursAudio outperforms all baselines supporting text input on recall (how many sound effects described in the text caption are generated),
which is the main metric we concern for the text-audio alignment axis. It also outperforms most baselines on precision and is on par with ElevenLabs on the real videos.
It should also be noted that CLAP score between TV2A are similar, which does not correlate strongly with the relative ranking.
However, CLAP is much higher for TV2A models compared to V2A models,
which shows its discriminative power when the delta is large enough,
as we expect V2A models to generate audio that has much worse alignment with text not they are not conditioned on.

For audio quality, we find reasonable correlation between AQual and subjective pairwise comparison.
On pairwise subjective tests \OursAudio outperforms all baselines,
and PikaLabs and ElevenLabs have the smallest gaps.
On the objective metrics, \OursAudio also has the highest score, followed by those two models.

On video-SFX alignment, the correlation between objective and subjective is much weaker.
We first note that Seeing\&Hearing directly use ImageBind for classifier-guidance which maximizes that.
Therefore, both the V2A and TV2A variant show IB score on par with \OursAudio, despite that the subjective tests reveal that the video-SFX alignment is much worse compared to \OursAudio.
Next, we observe that \OursAudio achieves the higher IB score compared to other baselines, and also outperforms them on subjective evaluations.
However the ranking for the baselines in terms of relative performance to \OursAudio on subjective tests does not correlate highly with IB ranking.

Objective results on \OursAudioBench for sound effect generation are shown in \cref{tab:app_eval_audio_main_ss_sfx_audiobench}. We note that similar trends are observed, where \OursAudio leads in AQual and IB (except when compared to Seeing\&Hearing which optimizes IB in inference), and on par with baselines on CLAP.

\subsubsection{Control audio quality through text prompts}\label{app:qual_control}
\cref{fig:audio_ablation_audio_quality_sfx} and \cref{fig:audio_ablation_audio_quality_sfx_music} present examples for audio quality control via text prompts for sound effect generation and joint music and sound effect generation respectively.
When conditioned on prompts indicating low quality (\ie, quality score of 5.0), \OursAudio generates natural but low quality audio that contains for example wind noises or music with broken bass.

\begin{table}[h]
    \centering
    \captionsetup{type=figure}
    \setlength{\tabcolsep}{1pt}
    \adjustbox{max width=0.95\textwidth}{%
    \centering
    \begin{tabular}{cccccc}
        \rowcolor{blue300}
        \multicolumn{6}{c}{\textbf{Ablation: SFX audio generation with different audio quality scores prompts (10s, \OursVideo)}} \\
        \includegraphics[width=0.17\linewidth]{figures/audio/teaser/ablation_aes_sfx/sample_198_aes_5/out-042.jpg}&
        \includegraphics[width=0.17\linewidth]{figures/audio/teaser/ablation_aes_sfx/sample_198_aes_5/out-085.jpg}&
        \includegraphics[width=0.17\linewidth]{figures/audio/teaser/ablation_aes_sfx/sample_198_aes_5/out-128.jpg}&
        \includegraphics[width=0.17\linewidth]{figures/audio/teaser/ablation_aes_sfx/sample_198_aes_5/out-170.jpg}&
        \includegraphics[width=0.17\linewidth]{figures/audio/teaser/ablation_aes_sfx/sample_198_aes_5/out-213.jpg}&
        \includegraphics[width=0.17\linewidth]{figures/audio/teaser/ablation_aes_sfx/sample_198_aes_5/out-256.jpg}\\
        \multicolumn{6}{c}{{(a) Audio quality score = 5}} \\
        \multicolumn{6}{c}{
            \includegraphics[width=\linewidth, height=0.1\linewidth]
            {figures/audio/teaser/ablation_aes_sfx/sample_198_aes_5/spec.jpg}
        } \\
        \multicolumn{6}{c}{{(b) Audio quality score = 6}} \\
        \multicolumn{6}{c}{
            \includegraphics[width=\linewidth, height=0.1\linewidth]
            {figures/audio/teaser/ablation_aes_sfx/sample_198_aes_6/spec.jpg}
        } \\
        \multicolumn{6}{c}{{(c) Audio quality score = 7}} \\
        \multicolumn{6}{c}{
            \includegraphics[width=\linewidth, height=0.1\linewidth]
            {figures/audio/teaser/ablation_aes_sfx/sample_198_aes_7/spec.jpg}
        } \\
        \multicolumn{6}{c}{{(d) Audio quality score = 8}} \\
        \multicolumn{6}{c}{
            \includegraphics[width=\linewidth, height=0.1\linewidth]
            {figures/audio/teaser/ablation_aes_sfx/sample_198_aes_8/spec.jpg}
        } \\
        \multicolumn{6}{c}{{(e) Audio quality score = 9}} \\
        \multicolumn{6}{c}{
            \includegraphics[width=\linewidth, height=0.1\linewidth]
            {figures/audio/teaser/ablation_aes_sfx/sample_198_aes_9/spec.jpg}
        } \\
    \end{tabular}}
\vspace{-3mm}
\caption{\textbf{Examples of generated audio with different audio quality score in the text prompt with \OursAudio.}
For these samples, higher audio quality scores tends to produced audio without wind noises compared to lower audio quality scores.
Videos in this Figure found at \url{https://go.fb.me/MovieGen-Figure40}.
}
\label{fig:audio_ablation_audio_quality_sfx}
\end{table}


\begin{table}[h]
    \centering
    \captionsetup{type=figure}
    \setlength{\tabcolsep}{1pt}
    \adjustbox{max width=0.95\textwidth}{%
    \centering
    \begin{tabular}{cccccc}
        \rowcolor{blue300}
        \multicolumn{6}{c}{\textbf{Ablation: SFX+music audio generation with different audio quality scores prompts (10s, \OursVideo)}} \\
        \includegraphics[width=0.17\linewidth]{figures/audio/teaser/ablation_aes_sfx_music/sample_231_aes_5/out-042.jpg}& 
        \includegraphics[width=0.17\linewidth]{figures/audio/teaser/ablation_aes_sfx_music/sample_231_aes_5/out-085.jpg}& 
        \includegraphics[width=0.17\linewidth]{figures/audio/teaser/ablation_aes_sfx_music/sample_231_aes_5/out-128.jpg}& 
        \includegraphics[width=0.17\linewidth]{figures/audio/teaser/ablation_aes_sfx_music/sample_231_aes_5/out-170.jpg}& 
        \includegraphics[width=0.17\linewidth]{figures/audio/teaser/ablation_aes_sfx_music/sample_231_aes_5/out-213.jpg}& 
        \includegraphics[width=0.17\linewidth]{figures/audio/teaser/ablation_aes_sfx_music/sample_231_aes_5/out-256.jpg}\\
        \multicolumn{6}{c}{{(a) Audio quality score = 5}} \\
        \multicolumn{6}{c}{
            \includegraphics[width=\linewidth, height=0.1\linewidth]
            {figures/audio/teaser/ablation_aes_sfx_music/sample_231_aes_5/spec.jpg}
        } \\
        \multicolumn{6}{c}{{(b) Audio quality score = 6}} \\
        \multicolumn{6}{c}{
            \includegraphics[width=\linewidth, height=0.1\linewidth]
            {figures/audio/teaser/ablation_aes_sfx_music/sample_231_aes_6/spec.jpg}
        } \\
        \multicolumn{6}{c}{{(c) Audio quality score = 7}} \\
        \multicolumn{6}{c}{
            \includegraphics[width=\linewidth, height=0.1\linewidth]
            {figures/audio/teaser/ablation_aes_sfx_music/sample_231_aes_7/spec.jpg}
        } \\
        \multicolumn{6}{c}{{(d) Audio quality score = 8}} \\
        \multicolumn{6}{c}{
            \includegraphics[width=\linewidth, height=0.1\linewidth]
            {figures/audio/teaser/ablation_aes_sfx_music/sample_231_aes_8/spec.jpg}
        } \\
        \multicolumn{6}{c}{{(e) Audio quality score = 9}} \\
        \multicolumn{6}{c}{
            \includegraphics[width=\linewidth, height=0.1\linewidth]
            {figures/audio/teaser/ablation_aes_sfx_music/sample_231_aes_9/spec.jpg}
        } \\
    \end{tabular}}

    \vspace{-3mm}

    \caption{\textbf{Examples of generated audio with different audio quality score in the text prompt with \OursAudio.}
    \OursAudio controls the generated audio output quality using the input text prompts (refer to \ref{tab:audio_caption}). As we can observed in the spectrogram plot, lower quality scores contains some high frequency noises and by gradually increasing the quality scores, we got a cleaner audio spectrogram.
    Videos in this Figure found at \url{https://go.fb.me/MovieGen-Figure41}.
}
\label{fig:audio_ablation_audio_quality_sfx_music}
\end{table}



\subsubsection{Control music style through text prompts}\label{app:music_control}
\cref{fig:audio_ablation_text_music} shows an example of varying music captions for the same video.




\begin{table}[h]
    \centering
    \captionsetup{type=figure}
    \setlength{\tabcolsep}{1pt}
    \adjustbox{max width=0.95\textwidth}{%
    \centering
    \begin{tabular}{cccccc}
        \rowcolor{blue300}
        \multicolumn{6}{c}{\textbf{Ablation: SFX+music audio generation with different text prompts (10s, \OursVideo)}} \\
        \includegraphics[width=0.17\linewidth]{figures/audio/teaser/ablation_text_music/sample_9_music1/out-042.jpg}& 
        \includegraphics[width=0.17\linewidth]{figures/audio/teaser/ablation_text_music/sample_9_music1/out-085.jpg}& 
        \includegraphics[width=0.17\linewidth]{figures/audio/teaser/ablation_text_music/sample_9_music1/out-128.jpg}& 
        \includegraphics[width=0.17\linewidth]{figures/audio/teaser/ablation_text_music/sample_9_music1/out-170.jpg}& 
        \includegraphics[width=0.17\linewidth]{figures/audio/teaser/ablation_text_music/sample_9_music1/out-213.jpg}& 
        \includegraphics[width=0.17\linewidth]{figures/audio/teaser/ablation_text_music/sample_9_music1/out-256.jpg}\\
        \multicolumn{6}{c}{{(a) Music style caption: A beautiful and emotional classical orchestral track with a cinematic feel.}} \\
        \multicolumn{6}{c}{
            \includegraphics[width=\linewidth, height=0.05\linewidth]
            {figures/audio/teaser/ablation_text_music/sample_9_music1/spec.jpg}
        } \\
        \multicolumn{6}{c}{{(b) Music style caption: A dark and mysterious electronic track with a haunting melody.}} \\
        \multicolumn{6}{c}{
            \includegraphics[width=\linewidth, height=0.05\linewidth]
            {figures/audio/teaser/ablation_text_music/sample_13_music2/spec.jpg}
        } \\
        \multicolumn{6}{c}{{(c) Music style caption: A slow, emotional, and passionate country song.}} \\
        \multicolumn{6}{c}{
            \includegraphics[width=\linewidth, height=0.05\linewidth]
            {figures/audio/teaser/ablation_text_music/sample_16_music3/spec.jpg}
        } \\
    \end{tabular}}
\vspace{-3mm}
\caption{\textbf{Examples of generated audio with different music description in the text prompt with \OursAudio.}
Videos in this Figure found at \url{https://go.fb.me/MovieGen-Figure42}.
}
\label{fig:audio_ablation_text_music}
\end{table}


\subsubsection{Additional Audio Samples}\label{app:audio_samples}
\cref{fig:audio_more_sfx} and \cref{fig:audio_more_sfx_music} presents additional samples of sound effect generation and joint SFX + music generation, respectively.

\begin{table}[h]
    \centering
    \captionsetup{type=figure}
    \setlength{\tabcolsep}{1pt}
    \adjustbox{max width=0.95\textwidth}{%
    \centering
    \begin{tabular}{cccc}
        \rowcolor{blue300}
        \multicolumn{4}{c}{\textbf{Appendix: additional sound effect samples}} \\

        \includegraphics[width=0.25\linewidth]{figures/audio/teaser/app_more_sfx/mvga_99/out-128.jpg} &
        \includegraphics[width=0.25\linewidth]{figures/audio/teaser/app_more_sfx/mvga_77/out-128.jpg} &
        \includegraphics[width=0.25\linewidth]{figures/audio/teaser/app_more_sfx/mvga_443/out-128.jpg} &
        \includegraphics[width=0.25\linewidth]{figures/audio/teaser/app_more_sfx/mvga_536/out-128.jpg} \\
        
        \multicolumn{1}{c}{
            \includegraphics[width=0.25\linewidth, height=0.05\linewidth]{figures/audio/teaser/app_more_sfx/mvga_99/spec.jpg}
        }  &
        \multicolumn{1}{c}{
            \includegraphics[width=0.25\linewidth, height=0.05\linewidth]{figures/audio/teaser/app_more_sfx/mvga_77/spec.jpg}
        }  &
        \multicolumn{1}{c}{
            \includegraphics[width=0.25\linewidth, height=0.05\linewidth]{figures/audio/teaser/app_more_sfx/mvga_443/spec.jpg}
        }  &
        \multicolumn{1}{c}{
            \includegraphics[width=0.25\linewidth, height=0.05\linewidth]{figures/audio/teaser/app_more_sfx/mvga_536/spec.jpg}
        }  \\
        
        \includegraphics[width=0.25\linewidth]{figures/audio/teaser/app_more_sfx/mvga_168/out-128.jpg} &
        \includegraphics[width=0.25\linewidth]{figures/audio/teaser/app_more_sfx/mvga_403/out-128.jpg} &
        \includegraphics[width=0.25\linewidth]{figures/audio/teaser/app_more_sfx/mvga_428/out-128.jpg} &
        \includegraphics[width=0.25\linewidth]{figures/audio/teaser/app_more_sfx/mvga_78/out-128.jpg} \\
        
        \multicolumn{1}{c}{
            \includegraphics[width=0.25\linewidth, height=0.05\linewidth]{figures/audio/teaser/app_more_sfx/mvga_168/spec.jpg}
        }  &
        \multicolumn{1}{c}{
            \includegraphics[width=0.25\linewidth, height=0.05\linewidth]{figures/audio/teaser/app_more_sfx/mvga_403/spec.jpg}
        }  &
        \multicolumn{1}{c}{
            \includegraphics[width=0.25\linewidth, height=0.05\linewidth]{figures/audio/teaser/app_more_sfx/mvga_428/spec.jpg}
        }  &
        \multicolumn{1}{c}{
            \includegraphics[width=0.25\linewidth, height=0.05\linewidth]{figures/audio/teaser/app_more_sfx/mvga_78/spec.jpg}
        }  \\
        
        \includegraphics[width=0.25\linewidth]{figures/audio/teaser/app_more_sfx/mvga_79/out-128.jpg} &
        \includegraphics[width=0.25\linewidth]{figures/audio/teaser/app_more_sfx/mvga_53/out-128.jpg} &
        \includegraphics[width=0.25\linewidth]{figures/audio/teaser/app_more_sfx/mvga_85/out-128.jpg} &
        \includegraphics[width=0.25\linewidth]{figures/audio/teaser/app_more_sfx/mvga_56/out-128.jpg} \\
        
        \multicolumn{1}{c}{
            \includegraphics[width=0.25\linewidth, height=0.05\linewidth]{figures/audio/teaser/app_more_sfx/mvga_79/spec.jpg}
        }  &
        \multicolumn{1}{c}{
            \includegraphics[width=0.25\linewidth, height=0.05\linewidth]{figures/audio/teaser/app_more_sfx/mvga_53/spec.jpg}
        }  &
        \multicolumn{1}{c}{
            \includegraphics[width=0.25\linewidth, height=0.05\linewidth]{figures/audio/teaser/app_more_sfx/mvga_85/spec.jpg}
        }  &
        \multicolumn{1}{c}{
            \includegraphics[width=0.25\linewidth, height=0.05\linewidth]{figures/audio/teaser/app_more_sfx/mvga_56/spec.jpg}
        }  \\
        
        \includegraphics[width=0.25\linewidth]{figures/audio/teaser/app_more_sfx/mvga_474/out-128.jpg} &
        \includegraphics[width=0.25\linewidth]{figures/audio/teaser/app_more_sfx/mvga_110/out-128.jpg} \\
        
        \multicolumn{1}{c}{
            \includegraphics[width=0.25\linewidth, height=0.05\linewidth]{figures/audio/teaser/app_more_sfx/mvga_474/spec.jpg}
        }  &
        \multicolumn{1}{c}{
            \includegraphics[width=0.25\linewidth, height=0.05\linewidth]{figures/audio/teaser/app_more_sfx/mvga_110/spec.jpg}
        }  \\
    \end{tabular}}
\vspace{-3mm}
\caption{\textbf{Additional \OursAudio sound effect samples.}
Videos in this Figure found at \url{https://go.fb.me/MovieGen-Figure43}.
}
\label{fig:audio_more_sfx}
\end{table}


\begin{table}[h]
    \centering
    \captionsetup{type=figure}
    \setlength{\tabcolsep}{1pt}
    \adjustbox{max width=0.95\textwidth}{%
    \centering
    \begin{tabular}{cccc}
        \rowcolor{blue300}
        \multicolumn{4}{c}{\textbf{Appendix: additional sound effect + music samples}} \\
        
        \includegraphics[width=0.25\linewidth]{figures/audio/teaser/app_more_sfx_music/mvga_88/out-128.jpg} &
        \includegraphics[width=0.25\linewidth]{figures/audio/teaser/app_more_sfx_music/mvga_290/out-128.jpg} &
        \includegraphics[width=0.25\linewidth]{figures/audio/teaser/app_more_sfx_music/mvga_286/out-128.jpg} &
        \includegraphics[width=0.25\linewidth]{figures/audio/teaser/app_more_sfx_music/mvga_169/out-128.jpg}\\
        
        \multicolumn{1}{c}{
            \includegraphics[width=0.25\linewidth, height=0.05\linewidth]{figures/audio/teaser/app_more_sfx_music/mvga_88/spec.jpg}
        }  &
        \multicolumn{1}{c}{
            \includegraphics[width=0.25\linewidth, height=0.05\linewidth]{figures/audio/teaser/app_more_sfx_music/mvga_290/spec.jpg}
        }  &
        \multicolumn{1}{c}{
            \includegraphics[width=0.25\linewidth, height=0.05\linewidth]{figures/audio/teaser/app_more_sfx_music/mvga_286/spec.jpg}
        }  &
        \multicolumn{1}{c}{
            \includegraphics[width=0.25\linewidth, height=0.05\linewidth]{figures/audio/teaser/app_more_sfx_music/mvga_169/spec.jpg}
        } \\
        
        \includegraphics[width=0.25\linewidth]{figures/audio/teaser/app_more_sfx_music/mvga_430/out-128.jpg} &
        \includegraphics[width=0.25\linewidth]{figures/audio/teaser/app_more_sfx_music/mvga_208/out-128.jpg} &
        \includegraphics[width=0.25\linewidth]{figures/audio/teaser/app_more_sfx_music/mvga_427/out-128.jpg} &
        \includegraphics[width=0.25\linewidth]{figures/audio/teaser/app_more_sfx_music/mvga_358/out-128.jpg}\\
        
        \multicolumn{1}{c}{
            \includegraphics[width=0.25\linewidth, height=0.05\linewidth]{figures/audio/teaser/app_more_sfx_music/mvga_430/spec.jpg}
        }  &
        \multicolumn{1}{c}{
            \includegraphics[width=0.25\linewidth, height=0.05\linewidth]{figures/audio/teaser/app_more_sfx_music/mvga_208/spec.jpg}
        }  &
        \multicolumn{1}{c}{
            \includegraphics[width=0.25\linewidth, height=0.05\linewidth]{figures/audio/teaser/app_more_sfx_music/mvga_427/spec.jpg}
        }  &
        \multicolumn{1}{c}{
            \includegraphics[width=0.25\linewidth, height=0.05\linewidth]{figures/audio/teaser/app_more_sfx_music/mvga_358/spec.jpg}
        }  \\
        
        \includegraphics[width=0.25\linewidth]{figures/audio/teaser/app_more_sfx_music/mvga_165/out-128.jpg} &
        \includegraphics[width=0.25\linewidth]{figures/audio/teaser/app_more_sfx_music/mvga_275/out-128.jpg} &
        \includegraphics[width=0.25\linewidth]{figures/audio/teaser/app_more_sfx_music/mvga_87/out-128.jpg} &
        \includegraphics[width=0.25\linewidth]{figures/audio/teaser/app_more_sfx_music/mvga_115/out-128.jpg}\\
        
        \multicolumn{1}{c}{
            \includegraphics[width=0.25\linewidth, height=0.05\linewidth]{figures/audio/teaser/app_more_sfx_music/mvga_165/spec.jpg}
        }  &
        \multicolumn{1}{c}{
            \includegraphics[width=0.25\linewidth, height=0.05\linewidth]{figures/audio/teaser/app_more_sfx_music/mvga_275/spec.jpg}
        }  &
        \multicolumn{1}{c}{
            \includegraphics[width=0.25\linewidth, height=0.05\linewidth]{figures/audio/teaser/app_more_sfx_music/mvga_87/spec.jpg}
        }  &
        \multicolumn{1}{c}{
            \includegraphics[width=0.25\linewidth, height=0.05\linewidth]{figures/audio/teaser/app_more_sfx_music/mvga_115/spec.jpg}
        } \\
        
        \includegraphics[width=0.25\linewidth]{figures/audio/teaser/app_more_sfx_music/mvga_514/out-128.jpg} &
        \includegraphics[width=0.25\linewidth]{figures/audio/teaser/app_more_sfx_music/mvga_339/out-128.jpg} &
        \includegraphics[width=0.25\linewidth]{figures/audio/teaser/app_more_sfx_music/mvga_265/out-128.jpg}\\
        
        \multicolumn{1}{c}{
            \includegraphics[width=0.25\linewidth, height=0.05\linewidth]{figures/audio/teaser/app_more_sfx_music/mvga_514/spec.jpg}
        }  &
        \multicolumn{1}{c}{
            \includegraphics[width=0.25\linewidth, height=0.05\linewidth]{figures/audio/teaser/app_more_sfx_music/mvga_339/spec.jpg}
        }  &
        \multicolumn{1}{c}{
            \includegraphics[width=0.25\linewidth, height=0.05\linewidth]{figures/audio/teaser/app_more_sfx_music/mvga_265/spec.jpg}
        } \\

    \end{tabular}}
\vspace{-3mm}
\caption{\textbf{Additional \OursAudio sound effect and music samples.}
Videos in this Figure found at \url{https://go.fb.me/MovieGen-Figure44}.
}
\label{fig:audio_more_sfx_music}
\end{table}

\clearpage
\subsection{Additional Results from \OursVideo}
Here, we include some further example generations from \OursVideo in~\cref{fig:t2v_qual_ours_figure_app1,fig:t2v_qual_ours_figure_app2}


\begin{center}
    \centering
    \captionsetup{type=figure}
    \setlength{\tabcolsep}{1pt}
\adjustbox{max width=\textwidth}{%
\centering
\begin{tabular}{cccc}
    \multicolumn{4}{c}{\textit{Prompt}: A person harvesting clouds from a field, placing them in a basket.} \\
    \includegraphics[width=0.25\linewidth]{figures/T2V_qualitative_appendix1/videos/clouds/frames/out-001.png} &
    \includegraphics[width=0.25\linewidth]{figures/T2V_qualitative_appendix1/videos/clouds/frames/out-061.png} &
    \includegraphics[width=0.25\linewidth]{figures/T2V_qualitative_appendix1/videos/clouds/frames/out-073.png} &
    \includegraphics[width=0.25\linewidth]{figures/T2V_qualitative_appendix1/videos/clouds/frames/out-255.png} \\

    \multicolumn{4}{c}{\textit{Prompt}: cinematic trailer for a group of samoyed puppies learning to become chefs.} \\
    \includegraphics[width=0.25\linewidth]{figures/T2V_qualitative_appendix1/videos/dogs/frames/out-001.png} &
    \includegraphics[width=0.25\linewidth]{figures/T2V_qualitative_appendix1/videos/dogs/frames/out-051.png} &
    \includegraphics[width=0.25\linewidth]{figures/T2V_qualitative_appendix1/videos/dogs/frames/out-101.png} &
    \includegraphics[width=0.25\linewidth]{figures/T2V_qualitative_appendix1/videos/dogs/frames/out-151.png} \\

    \multicolumn{4}{c}{\textit{Prompt}: A monk meditating in a temple carved into the cliffs of Bhutan.} \\
    \includegraphics[width=0.25\linewidth]{figures/T2V_qualitative_appendix1/videos/monk/frames/out-001.png} &
    \includegraphics[width=0.25\linewidth]{figures/T2V_qualitative_appendix1/videos/monk/frames/out-051.png} &
    \includegraphics[width=0.25\linewidth]{figures/T2V_qualitative_appendix1/videos/monk/frames/out-151.png} &
    \includegraphics[width=0.25\linewidth]{figures/T2V_qualitative_appendix1/videos/monk/frames/out-255.png} \\

    \multicolumn{4}{c}{\textit{Prompt}: Rivers of lava flow through a landscape of ice and snow, with steam rising into the air.} \\
    \includegraphics[width=0.25\linewidth]{figures/T2V_qualitative_appendix1/videos/lava/frames/out-001.png} &
    \includegraphics[width=0.25\linewidth]{figures/T2V_qualitative_appendix1/videos/lava/frames/out-101.png} &
    \includegraphics[width=0.25\linewidth]{figures/T2V_qualitative_appendix1/videos/lava/frames/out-151.png} &
    \includegraphics[width=0.25\linewidth]{figures/T2V_qualitative_appendix1/videos/lava/frames/out-255.png} \\

    \multicolumn{4}{c}{\textit{Prompt}: A young explorer who discovers a cave filled with glowing crystals.} \\
    \includegraphics[width=0.25\linewidth]{figures/T2V_qualitative_appendix1/videos/crystal/frames/out-001.png} &
    \includegraphics[width=0.25\linewidth]{figures/T2V_qualitative_appendix1/videos/crystal/frames/out-051.png} &
    \includegraphics[width=0.25\linewidth]{figures/T2V_qualitative_appendix1/videos/crystal/frames/out-151.png} &
    \includegraphics[width=0.25\linewidth]{figures/T2V_qualitative_appendix1/videos/crystal/frames/out-255.png} \\

    \multicolumn{4}{c}{\textit{Prompt}: A potter crafting ceramics using volcanic ash in Hawaii.} \\
    \includegraphics[width=0.25\linewidth]{figures/T2V_qualitative_appendix1/videos/pottery/frames/out-001.png} &
    \includegraphics[width=0.25\linewidth]{figures/T2V_qualitative_appendix1/videos/pottery/frames/out-051.png} &
    \includegraphics[width=0.25\linewidth]{figures/T2V_qualitative_appendix1/videos/pottery/frames/out-171.png} &
    \includegraphics[width=0.25\linewidth]{figures/T2V_qualitative_appendix1/videos/pottery/frames/out-255.png} \\


\end{tabular}}

    \vspace{-3mm}
    \caption{\textbf{Generated videos from \OursVideo.} 
    Videos in this Figure found at \url{https://go.fb.me/MovieGen-Figure45}.}
    \label{fig:t2v_qual_ours_figure_app1}
\end{center}%



\begin{center}
    \centering
    \captionsetup{type=figure}
    \setlength{\tabcolsep}{1pt}
\adjustbox{max width=\textwidth}{%
\centering
\begin{tabular}{cccc}
    \multicolumn{4}{c}{\textit{Prompt}: A dense jungle pathway is illuminated by oversized, bioluminescent mushrooms.} \\
    \includegraphics[width=0.25\linewidth]{figures/T2V_qualitative_appendix2/videos/jungle/frames/out-001.png} &
    \includegraphics[width=0.25\linewidth]{figures/T2V_qualitative_appendix2/videos/jungle/frames/out-051.png} &
    \includegraphics[width=0.25\linewidth]{figures/T2V_qualitative_appendix2/videos/jungle/frames/out-151.png} &
    \includegraphics[width=0.25\linewidth]{figures/T2V_qualitative_appendix2/videos/jungle/frames/out-201.png} \\

    \multicolumn{4}{c}{\textit{Prompt}: A turtle in a racing suit, riding a skateboard down a steep hill.} \\
    \includegraphics[width=0.25\linewidth]{figures/T2V_qualitative_appendix2/videos/turtle/frames/out-101.png} &
    \includegraphics[width=0.25\linewidth]{figures/T2V_qualitative_appendix2/videos/turtle/frames/out-151.png} &
    \includegraphics[width=0.25\linewidth]{figures/T2V_qualitative_appendix2/videos/turtle/frames/out-201.png} &
    \includegraphics[width=0.25\linewidth]{figures/T2V_qualitative_appendix2/videos/turtle/frames/out-255.png} \\

    \multicolumn{4}{c}{\textit{Prompt}:a woman eating ice scream.} \\
    \includegraphics[width=0.25\linewidth]{figures/T2V_qualitative_appendix2/videos/ice_cream/frames/out-001.png} &
    \includegraphics[width=0.25\linewidth]{figures/T2V_qualitative_appendix2/videos/ice_cream/frames/out-051.png} &
    \includegraphics[width=0.25\linewidth]{figures/T2V_qualitative_appendix2/videos/ice_cream/frames/out-151.png} &
    \includegraphics[width=0.25\linewidth]{figures/T2V_qualitative_appendix2/videos/ice_cream/frames/out-201.png} \\

    \multicolumn{4}{c}{\textit{Prompt}: Sailboat sailing through the crystal-clear waters of Bora Bora. Camera aerial shot.} \\
    \includegraphics[width=0.25\linewidth]{figures/T2V_qualitative_appendix2/videos/boat/frames/out-101.png} &
    \includegraphics[width=0.25\linewidth]{figures/T2V_qualitative_appendix2/videos/boat/frames/out-151.png} &
    \includegraphics[width=0.25\linewidth]{figures/T2V_qualitative_appendix2/videos/boat/frames/out-201.png} &
    \includegraphics[width=0.25\linewidth]{figures/T2V_qualitative_appendix2/videos/boat/frames/out-255.png} \\

    \multicolumn{4}{c}{\textit{Prompt}: A young detective who solves the case of the glowing plants.} \\
    \includegraphics[width=0.25\linewidth]{figures/T2V_qualitative_appendix2/videos/doctor/frames/out-001.png} &
    \includegraphics[width=0.25\linewidth]{figures/T2V_qualitative_appendix2/videos/doctor/frames/out-051.png} &
    \includegraphics[width=0.25\linewidth]{figures/T2V_qualitative_appendix2/videos/doctor/frames/out-151.png} &
    \includegraphics[width=0.25\linewidth]{figures/T2V_qualitative_appendix2/videos/doctor/frames/out-201.png} \\

    \multicolumn{4}{c}{\textit{Prompt}: Giant panda riding a bike through the streets of Beijing. Camera tracking shot.} \\
    \includegraphics[width=0.25\linewidth]{figures/T2V_qualitative_appendix2/videos/panda/frames/out-001.png} &
    \includegraphics[width=0.25\linewidth]{figures/T2V_qualitative_appendix2/videos/panda/frames/out-101.png} &
    \includegraphics[width=0.25\linewidth]{figures/T2V_qualitative_appendix2/videos/panda/frames/out-201.png} &
    \includegraphics[width=0.25\linewidth]{figures/T2V_qualitative_appendix2/videos/panda/frames/out-255.png} \\


\end{tabular}}

    \vspace{-3mm}
    \caption{\textbf{Generated videos from \OursVideo.} 
    Videos in this Figure found at \url{https://go.fb.me/MovieGen-Figure46}.}
    \label{fig:t2v_qual_ours_figure_app2}
\end{center}%

